\section{Detección de Objetos mediante Aprendizaje Profundo en Entornos Urbanos}

La automatización de procesos de supervisión vehicular en infraestructura urbana
requiere la localización y clasificación precisa de las instancias presentes en
el flujo de video. En el control de accesos a parqueaderos, los algoritmos de
detección de objetos constituyen el núcleo del sistema de percepción, siendo
responsables de identificar vehículos completos (sedanes, camionetas, minibuses)
y aislar subregiones de interés como las placas de matrícula.

\subsection{Formulación de la detección: Cajas delimitadoras y Funciones de Pérdida}

Como formulan Redmon et al. \cite{redmon2016yolo}, la detección de objetos
reformula la clasificación tradicional al predecir simultáneamente la categoría
del objeto y un vector continuo de coordenadas espaciales que define una caja
delimitadora (\textit{bounding box}). Para cada objeto identificado en el
parqueadero, la red predice una tupla $B = (x, y, w, h)$, donde $(x, y)$
representan las coordenadas del centro de la caja ajustadas dinámicamente a la
celda de la Rejilla, mientras que $w$ y $h$ denotan el ancho y el alto
normalizados con respecto a las dimensiones totales de la imagen. Este principio
permite encerrar desde la silueta completa de un vehículo ingresando al recinto
hasta aislar con alta resolución espacial la región rectangular de la matrícula.

Para medir cuantitativamente el grado de solapamiento entre la caja delimitadora
predicha por el modelo ($B_{\text{pred}}$) y la caja de referencia anotada
manualmente ($B_{\text{gt}}$), se utiliza la métrica de Intersección sobre Unión
(\textit{Intersection over Union}, IoU) \cite{rezatofighi2019giou}:

\begin{equation}
    \text{IoU}(B_{\text{pred}}, B_{\text{gt}})
    =
    \frac{|B_{\text{pred}} \cap B_{\text{gt}}|}{|B_{\text{pred}} \cup B_{\text{gt}}|}
\end{equation}

Dado que la métrica IoU convencional produce gradientes nulos cuando dos cajas
no se solapan ($\text{IoU} = 0$), las arquitecturas modernas optimizan la
localización mediante funciones de pérdida geométricas continuas como
Generalized IoU ($\mathcal{L}_{\text{GIoU}}$) y Complete IoU
($\mathcal{L}_{\text{CIoU}}$) \cite{zheng2020diou}, las cuales incorporan
penalizaciones por la distancia entre centros de las cajas y la discrepancia en
su relación de aspecto:

\begin{equation}
    \mathcal{L}_{\text{CIoU}}
    =
    1 - \text{IoU} + \frac{\rho^2(b, b^{\text{gt}})}{c^2} + \alpha v
\end{equation}

donde $\rho(b, b^{\text{gt}})$ representa la distancia euclídea entre los puntos
centrales de la caja predicha y la real, $c$ es la diagonal del contenedor
convexo mínimo que abarca ambas cajas, y $v$ mide la consistencia en la relación
de aspecto geométrica de la placa o vehículo.

\subsection{Taxonomía de detectores: Modelos de una etapa para inferencia en tiempo real}

Históricamente, la literatura de visión artificial dividió los detectores
profundos en dos grandes familias: los detectores de dos etapas
(\textit{two-stage}) y los de una sola etapa (\textit{one-stage}). Como
describen Ren et al. \cite{ren2015faster}, los modelos de dos etapas (como
Faster R-CNN) emplean una Red de Propuesta de Regiones (\textit{Region Proposal
Network}, RPN) para generar candidatos de objetos, los cuales son posteriormente
recortados, vectorizados y clasificados por subredes secundarias. Aunque este
enfoque ofrece un nivel elevado de precisión, la evaluación secuencial de
cientos de propuestas regionales genera un cuello de botella computacional
insuperable para el procesamiento de video de alta tasa de fotogramas.

En el contexto del control vehicular en parqueaderos, la latencia de inferencia
es un factor crítico. Por ello, los detectores de dos etapas son descartados
frente a las arquitecturas de una sola etapa (\textit{one-stage}) como SSD
\cite{liu2016ssd} y la familia YOLO \cite{redmon2016yolo, wang2024yolov1to10}.
Estos modelos evalúan la escena completa en una única pasada convolucional
(\textit{single-pass}), realizando la regresión de coordenadas y la asignación
de probabilidades de clase de forma simultánea. Esta unificación del flujo de
datos reduce el tiempo de procesamiento a milisegundos por fotograma,
consolidando a las arquitecturas de una etapa como el estándar de facto para la
supervisión vehicular en tiempo real.

\subsection{Evolución arquitectónica de la familia YOLO}

La evolución de la familia YOLO ha estado marcada por un cambio continuo en la
representación de los datos y en la eficiencia computacional. Las primeras
generaciones (YOLOv2 a YOLOv5) \cite{redmon2016yolo} introdujeron el concepto de
cajas ancla (\textit{anchor boxes}), las cuales requerían predefinir manualmente
plantillas de forma y tamaño mediante algoritmos de agrupamiento ($k$-means)
sobre el conjunto de entrenamiento.

Sin embargo, como demuestran Tian et al. \cite{tian2019fcos} y Ge et al. en
YOLOX, la dependencia de anclas fijas introduce rigidez y una alta sensibilidad
a hiperparámetros. La transición hacia arquitecturas libres de anclas
(\textit{anchor-free}), adoptada en YOLOv8, YOLOv10, YOLO11 y YOLO26 
\cite{yolo11_ultralytics,sapkota2026yolo26, jocher2026yolo26}, eliminó las cajas de referencia
predefinidas. En su lugar, el detector predice directamente desde cada punto
sobre el mapa de características las cuatro distancias continuas $(l, t, r, b)$ 
hacia los bordes superior, inferior, izquierdo y derecho de la entidad visual.

Esta transición hacia el enfoque \textit{anchor-free} mejoró drásticamente la
capacidad de detección multiescala dentro de la misma escena de parqueadero. Al
no depender de dimensiones de ancla preajustadas, la red adapta dinámicamente
sus campos receptivos, permitiendo localizar con alta precisión objetos de
dimensiones dispares dentro del mismo encuadre: desde un camión de carga de gran
volumen hasta una pequeña placa de matrícula instalada en el parachoques de un
automóvil.

\subsection{Cuellos de botella en el Borde (Edge AI): NMS y DFL en oclusión vehicular}

El despliegue de sistemas inteligentes en puntos de acceso vehicular exige
ejecutar los modelos de visión directamente en dispositivos de borde
(\textit{Edge AI}) integrados en las casetas del parqueadero (como tarjetas
NVIDIA Jetson o procesadores embebidos), donde la memoria, la potencia de
cálculo y la disipación térmica son limitadas \cite{sapkota2026yolo26}. En este
escenario, dos componentes arquitectónicos tradicionales representan los
principales cuellos de botella computacionales:

\begin{enumerate}
    \item \textbf{Carga computacional de Distribution Focal Loss (DFL):} Módulos
        como DFL, derivados de la pérdida focal \cite{lin2017focal}, transforman
        la regresión de cada borde de la caja en una
        distribución de probabilidad discreta distribuida en $K$ contenedores
        ($K=16$). Evaluar funciones \textit{Softmax} y sumas ponderadas sobre
        miles de puntos espaciales incrementa significativamente las operaciones
        de punto flotante (FLOPs) y el tamaño del cabezal de detección. En
        dispositivos con recursos restringidos, DFL ralentiza el tiempo de
        respuesta y complica la exportación de modelos a formatos de ejecución
        cuantitizados como TensorRT o INT8 \cite{jocher2026yolo26}.
    
    \item \textbf{Latencia estocástica de Supresión No Máxima (NMS):} El
        algoritmo tradicional NMS ordena y elimina iterativamente las cajas con
        alto grado de solapamiento mediante un umbral rígido de IoU
        \cite{bodla2017softnms}. En accesos con tráfico denso, cuando los
        automóviles avanzan en fila muy cercanos o sufren oclusiones parciales,
        NMS presenta dos fallas críticas:

    \begin{itemize}
        \item \textit{Latencia variable:} El tiempo de ejecución en CPU/GPU
            oscila impredeciblemente en función de la cantidad de vehículos
            detectados en la escena.

        \item \textit{Falsos negativos por oclusión:} NMS corre el riesgo de
            suprimir erróneamente la caja delimitadora de un automóvil válido al
            interpretar que su alto nivel de solapamiento con la caja del
            vehículo delantero es un duplicado, provocando la pérdida de
            registros de acceso.
    \end{itemize}
\end{enumerate}

\subsection{Arquitecturas modernas (NMS-Free) y Detección de Objetos Pequeños}

Para eliminar los retrasos del post-procesamiento en entornos de acceso crítico,
las versiones más recientes como YOLOv10, YOLOv12 y YOLO26
\cite{wang2024yolov1to10, tian2025yolov12_arxiv, sapkota2026yolo26} introducen
un diseño de cabezal dual (\textit{dual-head}). Durante el entrenamiento, el
modelo utiliza una rama uno-a-muchos (\textit{one-to-many}) para proporcionar
una supervisión densa y acelerar la convergencia, y una rama uno-a-uno
(\textit{one-to-one}) que aprende a asignar exactamente un predictor positivo
por objeto. Durante la inferencia en producción, se utiliza exclusivamente la
rama uno-a-uno, logrando una decodificación directa \textit{End-to-End NMS-Free}
que erradica la dependencia de NMS y garantiza una latencia constante. Además,
YOLO26 elimina por completo el módulo DFL, aligerando el cabezal de regresión y
reduciendo en más del 20\% los FLOPs requeridos \cite{jocher2026yolo26}.

Por otro lado, la localización de matrículas lejanas representa el reto de
detectar ``objetos pequeños'' (\textit{small targets}), los cuales ocupan apenas
una mínima fracción de píxeles en la captura de alta resolución. En los métodos
de asignación tradicionales (\textit{Task-Aligned Learning}, TAL), las
operaciones de submuestreo (\textit{downsampling}) reducen la resolución
espacial de los mapas de características profundos de tal forma que los objetos
muy pequeños no logran contener ningún punto de ancla en su interior, recibiendo
cero asignaciones positivas y anulando el gradiente de aprendizaje para dicha
muestra \cite{jocher2026yolo26}.

Para resolver esta degradación, modelos como YOLOv13 y YOLO26 incorporan la
técnica de Asignación de Etiquetas Adaptativa para Objetos Pequeños
(\textit{Small-Target-Aware Label Assignment}, STAL) \cite{lei2025yolov13,
sapkota2026yolo26}. STAL desacopla la geometría utilizada durante el filtrado de
candidatos de la superficie real de la caja, garantizando de forma matemática la
generación de muestras positivas para objetos cuya superficie sea inferior al
salto del mapa de características (\textit{stride}). Gracias a esta innovación,
el detector mantiene una supervisión continua sobre las matrículas distantes, 
resolviendo el problema crítico de la pérdida de detección de patentes lejanas
en las barreras de ingreso.
